% Part: proof-theory
% Chapter: cut-elimination
% Section: interpolation

\documentclass[../../../include/open-logic-section]{subfiles}

\begin{document}

\olfileid{pt}{cut}{itp}

\olsection{માએહારાનું ઉપપ્રમેય અને ક્રેગનું અંતર્વેશન પ્રમેય}

વિલિયમ ક્રેગનું અંતર્વેશન પ્રમેય કહે છે:
જો $!A \Entails !B$, તો એવું
!!a{sentence}~$!C$ (\emph{અંતર્વેશક}) છે કે
$!A \Entails !C$ અને $!C \Entails !B$.
મહત્ત્વની વાત એ છે કે અંતર્વેશક~$!C$
એવો પસંદ કરી શકાય કે આ~$!C$માં આવતાં
બધાં !!{constant}s અને !!{predicate}s
~$!A$ તથા~$!B$ બંનેમાં આવે.
(અહીં કોઈ !!{function}s સામેલ નથી એમ ધારીએ છીએ.)
આ મર્યાદા ન હોય તો~$!A$ અને~$!B$
પોતે જ સરળતાથી અંતર્વેશક બની જાય.

અંતર્વેશન પ્રમેય શાસ્ત્રીય પ્રથમ-ક્રમ
તર્ક માટે સાચું છે. તેને શોધાયેલો
પ્રથમ-ક્રમ તર્કનો ``છેલ્લો મહત્ત્વનો ગુણધર્મ''
પણ કહેવામાં આવ્યો છે.
નિદર્શ સિદ્ધાંત અને સ્વયંસંચાલિત
તર્કવિચારમાં તેના મહત્ત્વના ઉપયોગો છે.
તેની મૂળ પ્રેરણા અને ઉપયોગ વિજ્ઞાનના
તત્ત્વજ્ઞાનમાં હતાં: કોઈ સિદ્ધાંતને
કયા મૂળભૂત ખ્યાલોથી સ્વયંસિદ્ધ કરી શકાય
કે ન શકાય, તે તપાસવામાં તે મદદ કરે છે.
તેમાંથી બેથનું વ્યાખ્યેયતા પ્રમેય પણ
ફલિત થાય છે: સિદ્ધાંત કોઈ ખ્યાલને
ગૂઢ રીતે વ્યાખ્યાયિત કરી શકે,
તો સ્પષ્ટ રીતે પણ કરી શકે છે.

ગણિતીય તર્કનાં અનેક પરિણામોની જેમ
અંતર્વેશન પ્રમેયને નિદર્શ સિદ્ધાંતની
અને સાબિતી સિદ્ધાંતની, બંને પદ્ધતિથી
સાબિત કરી શકાય છે.
સાબિતી-સૈદ્ધાંતિક પદ્ધતિનો લાભ એ છે કે
તે અંતર્વેશકનું અસ્તિત્વ જ નહીં,
પરંતુ $!A \Entails !B$ની
!!a{proof}માંથી તેને ગણતો કલનવિધિ
પણ આપે છે; દાખલા તરીકે,
~\Log{G3c}માં
$!A \Sequent !B$ની !!a{proof}માંથી.

$\Lang L(!A)$ એ~$!A$માંનાં
!!{constant}s અને !!{predicate}sનો ગણ
અને
$\Lang L(\Gamma) = \bigcup \Setabs{\Lang L(!A)}{!A \in
\Gamma}$ માનો. આ સમગ્ર વિભાગમાં
!!{function}s વિનાનાં !!{formula}s
લેવાં છે.

\begin{thm}[ક્રેગનું અંતર્વેશન પ્રમેય]\ollabel{thm:interpolation}
~$!A$ની !!{language}ને~$\Lang L(!A)$
અને~$!B$ની ભાષાને~$\Lang L(!B)$ કહો.
જો $\Log{G3c} \Proves !A \Sequent !B$,
તો એવી !!a{formula}~$!C$ છે
જેની !!{language} માટે
$\Lang L(!C) \subseteq \Lang
L(!A) \cap \Lang L(!B)$ અને
$\Log{G3c} \Proves !A \Sequent !C$
તેમજ
$\Log{G3c} \Proves !C \Sequent !B$.
\end{thm}

~\Log{G3c}ની !!{proof}sની રચનાનો
ઉપયોગ કરવા અંતર્વેશન પ્રમેયનું
સામાન્યીકરણ સાબિત કરીએ છીએ.
તે માટે સિક્વન્ટના પૂર્વાંગ અને
ઉત્તરાંગને બે બે ભાગમાં વહેંચેલા
વિચારો. $\maeh{\Gamma_1}{\Gamma_2} \Sequent
\maeh{\Delta_1}{\Delta_2}$ લખીએ ત્યારે
$\Gamma_1$ અને $\Delta_1$ પહેલા ભાગમાં,
અને $\Gamma_2$ અને $\Delta_2$
બીજા ભાગમાં છે. પછી નીચેના ઉપપ્રમેયથી
અંતર્વેશન પ્રમેય તરત ફલિત થાય છે.

\begin{lem}[માએહારાનું ઉપપ્રમેય]\ollabel{lem:maehara}
જો $\Log{G3c}
\Proves \maeh{\Gamma_1}{\Gamma_2} \Sequent \maeh{\Delta_1}{\Delta_2}$,
તો એવી !!a{formula}~$!C$ છે
જેની !!{language} માટે
$\Lang L(!C)
\subseteq \Lang L(\Gamma_1,
\Delta_1) \cap \Lang L(\Gamma_2, \Delta_2)$
અને
$\Log{G3c} \Proves \Gamma_1 \Sequent \Delta_1, !C$
તથા
$\Log{G3c} \Proves !C,
\Gamma_2 \Sequent \Delta_2$.
\end{lem}

\begin{proof}
  માનો કે $\pi \Proves \maeh{\Gamma_1}{\Gamma_2} \Sequent
  \maeh{\Delta_1}{\Delta_2}$.
  $n=\pheight{\pi}$ પર આગમનથી
  દાવો સાબિત કરીએ છીએ.

  જો $n=0$, તો
  $\maeh{\Gamma_1}{\Gamma_2} \Sequent
  \maeh{\Delta_1}{\Delta_2}$
  સ્વયંસિદ્ધ છે અને કોઈ પરમાણ્વીય
  !!{formula}~$!A$ એ
  $\Gamma_1, \Gamma_2$ તથા
  $\Delta_1, \Delta_2$ બંનેમાં આવે છે.
  ડાબે~$!A$ પહેલા કે બીજા ભાગમાં,
  અને જમણે પહેલા કે બીજા ભાગમાં છે,
  તેના ચાર કિસ્સા છે.
  \begin{enumerate}
    \item $!A \in \Gamma_1$ અને $!A \in \Delta_1$.
    એવું~$!C$ જોઈએ કે
    $\Gamma_1 \Sequent \Delta_1, !C$ અને
    $!C, \Gamma_2 \Sequent \Delta_2$
    બંને !!{provable} હોય.
    પહેલું તો~$!C$ ગમે તે લઈએ
    તોય સ્વયંસિદ્ધ છે, કારણ કે~$!A$
    ડાબે અને જમણે બંને છે.
    $!C, \Gamma_2 \Sequent \Delta_2$
    !!{provable} રહે તેની ખાતરી માટે
    $!C \ident \lfalse$ લેવું પડે.
    \item $!A \in \Gamma_1$ અને
    $!A \in \Delta_2$. ત્યારે
    $\Gamma_1 \Sequent \Delta_1, !A$
    અને $!A, \Gamma_2 \Sequent \Delta_2$
    બંને સ્વયંસિદ્ધ છે; તેથી
    $!C \ident !A$ લઈ શકાય.
    \item $!A \in \Gamma_2$ અને
    $!A \in \Delta_1$. ત્યારે
    $!A, \Gamma_1 \Sequent \Delta_1$
    અને $\Gamma_2 \Sequent \Delta_2, !A$
    સ્વયંસિદ્ધ હોય, પરંતુ તેમાં~$!A$
    ખોટી બાજુએ છે. આથી~$!A$
    અંતર્વેશક નથી. તેમ છતાં
    $\RightR{\lnot}$ અને
    $\LeftR{\lnot}$ વડે
    $\Gamma_1 \Sequent \Delta_1, \lnot!A$
    અને $\lnot !A, \Gamma_2 \Sequent \Delta_2$
    !!{prove} કરી શકાય.
    \item $!A \in \Gamma_2$ અને
    $!A \in \Delta_2$. કસરત.
  \end{enumerate}
  જો $\lfalse \in \Gamma_1$
  અથવા $\lfalse \in \Gamma_2$
  હોવાને કારણે
  $\maeh{\Gamma_1}{\Gamma_2} \Sequent
  \maeh{\Delta_1}{\Delta_2}$
  સ્વયંસિદ્ધ હોય, તો એ કિસ્સાઓ
  પણ કસરત તરીકે છોડીએ છીએ.

  જો $n>0$, તો~$\pi$નો અંત
  તાર્કિક અનુમાનથી થાય છે.
  કયો નિયમ વપરાયો અને મુખ્ય સૂત્ર
  કયા ભાગમાં છે, તેના કિસ્સા
  જુદા પાડવા પડે.
  દરેક કિસ્સામાં આગમન-પૂર્વધારણા
  વાપરી શકીએ: $\pheight{\pi_1} < n$
  ધરાવતી દરેક !!{proof}~$\pi_1$
  માટે, ~$\pi_1$ના અંતિમ સિક્વન્ટના
  બે ભાગ ગમે તે રીતે કર્યા હોય,
  ઉપપ્રમેય સાચું છે.
  હવે કેટલાક ચોક્કસ કિસ્સા જોઈએ.
  \begin{enumerate}
    \item $\pi$નો અંત~$\RightR{\lnot}$થી
    થાય છે અને મુખ્ય સૂત્ર~$\lnot !A$ છે.
    $\lnot !A$ પહેલા કે બીજા ભાગમાં છે
    તે પ્રમાણે~$\pi$નો અંત બેમાંથી
    એક રીતે થાય છે:
    \[
    \AxiomC{}
    \RightLabel{$\pi_1$}
    \Deduce$\maeh{!A, \Gamma_1}{\Gamma_2} \fCenter
      \maeh{\Delta_1}{\Delta_2}$
    \RightLabel{\RightR{\lnot}}
    \UnaryInf$\maeh{\Gamma_1}{\Gamma_2} \fCenter
      \maeh{\Delta_1, \lnot !A}{\Delta_2}$
    \DisplayProof
    \]
    અથવા
    \[
    \AxiomC{}
    \RightLabel{$\pi_1$}
    \Deduce$\maeh{\Gamma_1}{!A, \Gamma_2} \fCenter
      \maeh{\Delta_1}{\Delta_2}$
    \RightLabel{\RightR{\lnot}}
    \UnaryInf$\maeh{\Gamma_1}{\Gamma_2} \fCenter
      \maeh{\Delta_1}{\Delta_2, \lnot !A}$
    \DisplayProof
    \]
    ધ્યાન આપો: જો મુખ્ય !!{formula}
    ~$\lnot !A$ પહેલા ભાગમાં હોય,
    તો આધારવિધાનમાં સક્રિય
    !!{formula}~$!A$ પણ પહેલા
    ભાગમાં મૂક્યું છે. આ પસંદગી
    આપણી છે. સક્રિય સૂત્ર બીજા
    ભાગમાં મૂક્યું હોય તોય
    આગમન-પૂર્વધારણા લાગુ પડે,
    પરંતુ પછી અંતર્વેશક મળી
    શકતો નથી (કસરત: અજમાવી જુઓ).

    પહેલાં~$\lnot !A$ પહેલા ભાગમાં
    હોય તે સ્થિતિ લો.
    આગમન-પૂર્વધારણાથી~$\pi_1$ના અંતિમ
    સિક્વન્ટનો અંતર્વેશક, એટલે
    !!a{formula}~$!C_1$, એવો છે કે નીચેના
    બંને સિક્વન્ટ
    \begin{align*}
    !A, \Gamma_1 & \Sequent \Delta_1, !C_1 &&\text{અને} &
    !C_1, \Gamma_2 & \Sequent \Delta_2
    \intertext{સાબિત થઈ શકે છે. આપણને એવી !!a{formula}~$!C$ જોઈએ કે નીચેના બંને સિક્વન્ટ !!{provable} હોય:}
    \Gamma_1 & \Sequent \Delta_1, \lnot !A, !C &&\text{અને} &
    !C, \Gamma_2 & \Sequent \Delta_2
    \end{align*}
    સ્પષ્ટ રીતે~$!C \ident !C_1$ લઈ શકાય:
    જમણો સિક્વન્ટ !!{provable} છે
    એ આગમન-પૂર્વધારણાથી મળે છે,
    અને ડાબો સિક્વન્ટ
    આગમન-પૂર્વધારણાથી મળેલા સિક્વન્ટ
    પર~$\RightR{\lnot}$ લગાવતાં મળે છે.
    આગમન-પૂર્વધારણા એમ પણ કહે છે કે
    $\Lang L(!C_1) \subseteq \Lang L(!A, \Gamma_1,
    \Delta_1) \cap \Lang L(\Gamma_2, \Delta_2)$.
    કારણ કે $\Lang L(\lnot !A) = \Lang L(!A)$
    અને $\Lang L(!C) = \Lang L(!C_1)$,
    આપણને
    $\Lang L(!C_1) \subseteq
    \Lang L(\Gamma_1, \Delta_1, \lnot !A) \cap
    \Lang L(\Gamma_2, \Delta_2)$ મળે છે.

    બીજા કિસ્સામાં~$\lnot !A$ બીજા
    ભાગમાં હોવાથી સ્થિતિ થોડી જુદી છે.
    આધારવિધાનનું સક્રિય !!{formula}
    પણ બીજા ભાગમાં મૂકીને
    આગમન-પૂર્વધારણા આના પર લગાવીએ:
    \begin{align*}
    \Gamma_1 & \Sequent \Delta_1, !C_1 &&\text{અને} &
    !C_1, !A, \Gamma_2 & \Sequent \Delta_2
    \intertext{આ બંને સાબિત થઈ શકે છે. બીજાં સિક્વન્ટ પર ફરી \RightR{\lnot} લગાવતાં આની !!{proof}s મળે છે:}
    \Gamma_1 & \Sequent \Delta_1, !C_1 &&\text{અને} &
    !C_1, \Gamma_2 & \Sequent \Delta_2, \lnot !A
    \end{align*}
    $!C_1$ અંતર્વેશક હોય ત્યારે
    આપણને !!{provable} જોઈએ એવા
    બરાબર આ જ સિક્વન્ટ છે;
    ફરી~$!C \ident !C_1$ લઈએ.
    અગાઉની જેમ
    $\Lang L(!C_1) \subseteq
    \Lang L(\Gamma_1, \Delta_1) \cap
    \Lang L(\Gamma_2, \Delta_2, \lnot !A)$.
    \item $\pi$નો અંત~$\RightR{\lif}$થી
    થાય છે અને મુખ્ય
    !!{formula}~$!A \lif !B$ છે.
    $!A \lif !B$ પહેલા કે બીજા
    ભાગમાં છે તે પ્રમાણે બે કિસ્સા છે.
    !!{proof}~$\pi$નો અંત આ રીતે થાય છે:
    \[
    \AxiomC{}
    \RightLabel{$\pi_1$}
    \Deduce$\maeh{!A, \Gamma_1}{\Gamma_2} \fCenter
      \maeh{\Delta_1, !B}{\Delta_2}$
    \RightLabel{\RightR{\lif}}
    \UnaryInf$\maeh{\Gamma_1}{\Gamma_2} \fCenter
      \maeh{\Delta_1, !A \lif !B}{\Delta_2}$
    \DisplayProof
    \]
    અથવા
    \[
    \AxiomC{}
    \RightLabel{$\pi_1$}
    \Deduce$\maeh{\Gamma_1}{!A, \Gamma_2} \fCenter
      \maeh{\Delta_1}{\Delta_2, !B}$
    \RightLabel{\RightR{\lif}}
    \UnaryInf$\maeh{\Gamma_1}{\Gamma_2} \fCenter
      \maeh{\Delta_1}{\Delta_2, !A \lif !B}$
    \DisplayProof
    \]
    ફરી, આધારવિધાનનાં સક્રિય
    !!{formula}sને નિષ્કર્ષના મુખ્ય
    !!{formula}ના જ ભાગમાં મૂક્યાં છે.
    પહેલા કિસ્સામાં આગમન-પૂર્વધારણાથી
    આની !!{proof}s મળે છે:
    \begin{align*}
    !A, \Gamma_1 & \Sequent \Delta_1, !B, !C_1 &&\text{અને} & !C_1, \Gamma_2 & \Sequent \Delta_2
    \intertext{ડાબો સિક્વન્ટ \RightR{\lif}નું યોગ્ય આધારવિધાન છે. તેથી નીચેના સિક્વન્ટ !!{provable} છે:}
    \Gamma_1 & \Sequent \Delta_1, !A \lif !B, !C_1 &&\text{અને} & !C_1, \Gamma_2 & \Sequent \Delta_2
    \end{align*}
    આથી~$!C_1$ એ~$\pi$ના અંતિમ
    સિક્વન્ટનો પણ અંતર્વેશક છે;
    ફરી~$!C \ident !C_1$ લઈ શકાય.
    
    બીજો કિસ્સો એ જ રીતે થાય છે.
    \item $\pi$નો અંત~$\LeftR{\lif}$થી
    થાય છે અને મુખ્ય !!{formula}
    ~$!A \lif !B$ છે.
    જો $!A \lif !B$ પહેલા ભાગમાં હોય,
    તો !!{proof}~$\pi$નો અંત આ રીતે થાય છે:
    \[
    \AxiomC{}
    \RightLabel{$\pi_1$}
    \Deduce$\maeh{\Gamma_1}{\Gamma_2} \fCenter
      \maeh{\Delta_1, !A}{\Delta_2}$
    \AxiomC{}
    \RightLabel{$\pi_2$}
    \Deduce$\maeh{!B, \Gamma_1}{\Gamma_2} \fCenter
      \maeh{\Delta_1}{\Delta_2}$
    \BinaryInf$\maeh{!A \lif !B, \Gamma_1}{\Gamma_2} \fCenter
      \maeh{\Delta_1}{\Delta_2}$
    \DisplayProof
    \]
    આગમન-પૂર્વધારણા ચાર
    !!{provable} સિક્વન્ટ આપે છે:
    ડાબા આધારવિધાનના અંતર્વેશક~$!C_1$
    માટે બે અને બીજા આધારવિધાનના
    અંતર્વેશક~$!C_2$ માટે બે:
    \begin{align*}
    \Gamma_1 & \Sequent \Delta_1, !A, !C_1 &&\text{(1)} &
    !C_1, \Gamma_2 & \Sequent \Delta_2 && \text{(2)}\\
    !B, \Gamma_1 & \Sequent \Delta_1, !C_2 &&\text{(3)} &
    !C_2, \Gamma_2 & \Sequent \Delta_2 && \text{(4)}
    \end{align*}
    શિથિલીકરણ વડે (1)ના જમણે~$!C_2$
    અને (3)ના જમણે~$!C_1$ ઉમેરીએ,
    તો સિક્વન્ટ (1) અને (3)
    ~\LeftR{\lif}નાં આધારવિધાન બને:
    \[
    \AxiomC{}
    \Deduce$\Gamma_1 \fCenter \Delta_1, !A, !C_1, !C_2$
    \AxiomC{}
    \Deduce$!B, \Gamma_1 \fCenter \Delta_1, !C_1, !C_2$
    \RightLabel{\LeftR{\lif}}
    \BinaryInf$!A \lif !B, \Gamma_1 \fCenter \Delta_1, !C_1, !C_2$
    \DisplayProof
    \]
    ~$\RightR{\lor}$ લગાવતાં સિક્વન્ટ મળે:
    \[
    !A \lif !B, \Gamma_1 \fCenter \Delta_1, !C_1 \lor !C_2
    \]
    આથી આ કિસ્સામાં
    $!C_1 \lor !C_2$ અંતર્વેશક હોઈ શકે.
    ખરેખર, બાકી સિક્વન્ટ (2) અને~(4)
    વડે બીજું જરૂરી સિક્વન્ટ
    !!{prove} કરી શકાય:
    \[
    \AxiomC{}
    \Deduce$!C_1, \Gamma_2 \fCenter \Delta_2$
    \AxiomC{}
    \Deduce$!C_2, \Gamma_2 \fCenter \Delta_2$
    \RightLabel{\LeftR{\lor}}
    \BinaryInf$!C_1 \lor !C_2, \Gamma_2 \fCenter \Delta_2$
    \DisplayProof
    \]
    $!C_1 \lor !C_2$ યોગ્ય ભાષામાં છે
    તે ચકાસવાનું રહે છે.
    આગમન-પૂર્વધારણાથી મળે છે કે
    \begin{align*}
    \Lang L(!C_1) & \subseteq
    \Lang L(\Gamma_1, \Delta_1, !A) \cap \Lang L(\Gamma_2, \Delta_2) &\text{અને}\\
    \Lang L(!C_2) & \subseteq
    \Lang L(!B, \Gamma_1, \Delta_1) \cap \Lang L(\Gamma_2, \Delta_2)
    \intertext{કારણ કે}
    \Lang L(!A) & \subseteq \Lang L(!A \lif !B) &\text{અને}\\
    \Lang L(!B) & \subseteq \Lang L(!A \lif !B)
    \intertext{તેથી બંને માટે}
    \Lang L(!C_1) & \subseteq
    \Lang L(\Gamma_1, \Delta_1, !A \lif !B) \cap \Lang L(\Gamma_2, \Delta_2) &\text{અને}\\
    \Lang L(!C_2) & \subseteq
    \Lang L(\Gamma_1, \Delta_1, !A \lif !B) \cap \Lang L(\Gamma_2, \Delta_2)
    \intertext{અને તેથી $\Lang L (!C_1 \lor !C_2) = {}$}
    \Lang L(!C_1) \cup \Lang L(!C_2) & \subseteq
    \Lang L(\Gamma_1, \Delta_1, !A \lif !B) \cap \Lang L(\Gamma_2, \Delta_2),
    \end{align*}
    જે જરૂરી હતું તે જ છે.

    જો~$!A \lif !B$ બીજા ભાગમાં હોય,
    તો સ્થિતિ જુદી છે, પણ દલીલ સમાન છે.
    આગમન-પૂર્વધારણા ફરી ચાર
    !!{provable} સિક્વન્ટ આપે છે:
    \begin{align*}
    \Gamma_1 & \Sequent \Delta_1, !C_1 &&\text{(1)} &
    !C_1, \Gamma_2 & \Sequent \Delta_2, !A && \text{(2)}\\
    \Gamma_1 & \Sequent \Delta_1, !C_2 &&\text{(3)} &
    !C_2, !B, \Gamma_2 & \Sequent \Delta_2 && \text{(4)}
    \end{align*}
    હવે સિક્વન્ટ (2) અને (4),
    શિથિલીકરણથી ઉમેરેલાં કેટલાંક
    !!{formula}s સાથે,
    ~\LeftR{\lif}નાં આધારવિધાન બને:
    \[
    \AxiomC{}
    \Deduce$!C_1, !C_2, \Gamma_2 \fCenter \Delta_2, !A$
    \AxiomC{}
    \Deduce$!B, !C_1, !C_2, \Gamma_2 \fCenter \Delta_2$
    \RightLabel{\LeftR{\lif}}
    \BinaryInf$!A \lif !B, !C_1, !C_2, \Gamma_2 \fCenter \Delta_2$
    \DisplayProof
    \]
    આ વખતે પણ~$!C_1$ અને~$!C_2$માંથી
    બનેલું એક જ !!{formula} ધરાવતો
    સિક્વન્ટ જોઈએ, પણ તે પૂર્વાંગમાં
    જોઈએ (કારણ કે હવે બીજો ભાગ છે).
    ~$\LeftR{\land}$ લગાવતાં કામ થાય છે;
    $!C \ident (!C_1 \land !C_2)$
    અંતર્વેશક લેવો જોઈએ.
    બાકી સિક્વન્ટ (1) અને (3)
    વડે બીજા ભાગને અનુરૂપ સિક્વન્ટ
    !!{prove} કરી શકાય:
    \[
    \AxiomC{}
    \Deduce$\Gamma_1 \fCenter \Delta_1, !C_1$
    \AxiomC{}
    \Deduce$\Gamma_1 \fCenter \Delta_1, !C_2$
    \RightLabel{\RightR{\land}}
    \BinaryInf$\Gamma_1 \fCenter \Delta_1, !C_1 \land !C_2$
    \DisplayProof
    \]
    અગાઉની જેમ
    $\Lang L(!C_1 \land !C_2) \subseteq \Lang
    L(\Gamma_1, \Delta_1) \cap
    \Lang L(\Gamma_2, \Delta_2, !A \lif !B)$.

    \item $\pi$નો અંત~$\RightR{\lforall}$થી
    થાય છે અને મુખ્ય !!{formula}
    ~$\lforall[x][!A(x)]$ છે:
    \[
    \AxiomC{}
    \RightLabel{$\pi_1$}
    \Deduce$\maeh{\Gamma_1}{\Gamma_2} \fCenter
    \maeh{\Delta_1, !A(c)}{\Delta_2}$
    \RightLabel{\RightR{\lforall}}
    \UnaryInf$\maeh{\Gamma_1}{\Gamma_2} \fCenter
      \maeh{\Delta_1, \lforall[x][!A(x)]}{\Delta_2}$
    \DisplayProof
    \]
    અથવા
    \[
    \AxiomC{}
    \RightLabel{$\pi_1$}
    \Deduce$\maeh{\Gamma_1}{\Gamma_2} \fCenter
      \maeh{\Delta_1}{\Delta_2, !A(c)}$
    \RightLabel{\RightR{\lforall}}
    \UnaryInf$\maeh{\Gamma_1}{\Gamma_2} \fCenter
      \maeh{\Delta_1}{\Delta_2, \lforall[x][!A(x)]}$
    \DisplayProof
    \]
    પહેલા કિસ્સામાં આગમન-પૂર્વધારણાથી
    $\Gamma_1 \Sequent \Delta_1, !A(c), !C_1$
    અને $!C_1, \Gamma_2 \Sequent \Delta_2$ની
    !!{proof}s મળે છે. પહેલા સિક્વન્ટ પર
    $\RightR{\lforall}$ લગાવતાં
    $\Gamma_1 \Sequent \Delta_1,
    \lforall[x][!A(x)], !C_1$ મળે છે.
    (સ્વચલ-શરત પૂરી છે, કારણ કે~$\pi$ના
    અંતિમ \RightR{\lforall}માં તે
    પહેલેથી પૂરી હતી.)
    તેથી
    $\Lang L(!C_1) \subseteq
    \Lang L(\Gamma_1, \Delta_1, \lforall[x][!A(x)])
    \cap \Lang L(\Gamma_2, \Delta_2)$
    હોય તો~$!C_1$ અંતર્વેશક છે.
    આગમન-પૂર્વધારણા
    $\Lang L(!C_1) \subseteq
    \Lang L(\Gamma_1, \Delta_1, !A(c))
    \cap \Lang L(\Gamma_2, \Delta_2)$ આપે છે.
    એટલે~$c$ અંગે ચિંતા છે:
    શું~$c$ એ~$!C_1$માં આવી શકે?
    ના. જો~$c$ એ~$!C_1$માં આવે,
    તો $c \in \Lang L(\Gamma_1, \Delta_1, !A(c))$
    (જે સાચું છે) અને
    $c \in \Lang L(\Gamma_2, \Delta_2)$
    બંને થશે. પરંતુ ત્યારે~$c$
    એ~$\pi$ના $\RightR{\lforall}$
    અનુમાનના નિષ્કર્ષમાં આવશે,
    જે સ્વચલ-શરતનો ભંગ કરે છે.
    
    બીજો કિસ્સો સમાન છે.

    \item $\pi$નો અંત~$\RightR{\lexists}$થી
    થાય છે અને મુખ્ય
    !!{formula}~$\lexists[x][!A(x)]$ છે.
    ફરી બે કિસ્સા છે.
    $\lexists[x][!A(x)]$ પહેલા ભાગમાં
    હોય તે કિસ્સો જોઈએ;
    બીજો કસરત તરીકે છોડીએ છીએ.

    પહેલા કિસ્સામાં~$\pi$નો અંત
    આ રીતે થાય છે:
    \[
    \AxiomC{}
    \RightLabel{$\pi_1$}
    \Deduce$\maeh{\Gamma_1}{\Gamma_2} \fCenter
      \maeh{\Delta_1, \lexists[x][!A(x)], !A(t)}{\Delta_2}$
    \RightLabel{\RightR{\lexists}}
    \UnaryInf$\maeh{\Gamma_1}{\Gamma_2} \fCenter
      \maeh{\Delta_1, \lexists[x][!A(x)]}{\Delta_2}$
    \DisplayProof
    \]
    આગમન-પૂર્વધારણાથી
    $\Gamma_1 \Sequent
    \Delta_1, \lexists[x][!A(x)], !A(t), !C_1$
    અને $!C_1, \Gamma_2 \Sequent \Delta_2$ની
    !!{proof}s મળે છે.
    પહેલા સિક્વન્ટ પર
    $\RightR{\lexists}$ લગાવતાં
    $\Gamma_1 \Sequent \Delta_1,
    \lexists[x][!A(x)], !C_1$ મળે છે.
    તેથી
    $\Lang L(!C_1) \subseteq
    \Lang L(\Gamma_1, \Delta_1,
    \lexists[x][!A(x)])
    \cap \Lang L(\Gamma_2, \Delta_2)$
    હોય તો~$!C_1$ ફરી અંતર્વેશક છે.
    અહીં આગમન-પૂર્વધારણાથી
    $\Lang L(!C_1) \subseteq
    \Lang L(\Gamma_1, \Delta_1,
    \lexists[x][!A(x)], !A(t))
    \cap \Lang L(\Gamma_2, \Delta_2)$ મળે છે.
    યાદ કરો કે !!{function}s નથી.
    પદ~$t$ ચલ હોઈ શકે; એ
    ભાષામાં નવું અતાર્કિક ચિહ્ન ઉમેરતું નથી,
    એટલે એ કિસ્સામાં અંતર્વેશક સીધો ચાલે.
    અચલના કિસ્સામાં તે
    !!a{constant} છે, એટલે
    $t \ident b$.
    \sourcecorrection{OLMD-206}{મૂળ લખાણમાં !!{function}s ન હોવાથી $t$ માત્ર અચલ છે એવો દાવો છે; ચલ પણ પદ છે, પરંતુ તે અંતર્વેશકની અતાર્કિક શબ્દાવલી વધારતું નથી. નીચેની અચલ-વિભાગની દલીલ યથાવત્ છે.}
    જો $b \notin \Lang L(!C_1)$
    અથવા
    $b \in \Lang L(\Gamma_1, \Delta_1,
    \lexists[x][!A(x)]) \cap
    \Lang L(\Gamma_2, \Delta_2)$,
    તો ચિંતા નથી:~$!C_1$ને
    અંતર્વેશક લઈ શકાય.
    પરંતુ જો
    $b \in \Lang L(!C_1)$ અને
    $b \notin \Lang
    L(\Gamma_1, \Delta_1,
    \lexists[x][!A(x)])
    \cap \Lang L(\Gamma_2, \Delta_2)$ હોય તો?
    પહેલાં~$!C_1$માં~$b$ના
    દરેક દૃષ્ટાંતને !!a{variable}~$y$થી
    બદલી
    $!C_1 \ident \Subst{!C_1'}{b}{y}$
    લખી શકાય, જ્યાં~$!C_1'$માં~$b$
    આવતો નથી. વધુમાં, કાં તો
    $b \notin
    \Lang L(\Gamma_1, \Delta_1,
    \lexists[x][!A(x)])$
    અથવા
    $b \notin \Lang L(\Gamma_2, \Delta_2)$.
    અનુક્રમે
    $!C \ident \lforall[y][!C_1'(y)]$
    અથવા
    $!C \ident \lexists[y][!C_1'(y)]$
    લઈ શકાય. પહેલા કિસ્સામાં:
    \[
    \AxiomC{}
    \Deduce$\Gamma_1 \fCenter
      \Delta_1, \lexists[x][!A(x)], !A(b), !C_1'(b)$
    \RightLabel{\RightR{\lexists}}
    \UnaryInf$\Gamma_1 \fCenter \Delta_1, \lexists[x][!A(x)], !C_1'(b)$
    \RightLabel{\RightR{\lforall}}
    \UnaryInf$\Gamma_1 \fCenter
      \Delta_1, \lexists[x][!A(x)], \lforall[y][!C_1'(y)]$
    \DisplayProof
    \]
    અને
    \[
    \AxiomC{}
    \Deduce$!C_1'(b), \lforall[y][!C_1'(y)], \Gamma_2 \fCenter \Delta_2$
    \RightLabel{\LeftR{\lforall}}
    \UnaryInf$\lforall[y][!C_1'(y)], \Gamma_2 \fCenter \Delta_2$
    \DisplayProof
    \]
    ઉપરના સિક્વન્ટની !!{proof}s
    આગમન-પૂર્વધારણાથી મળે છે
    (બીજામાં પૂર્વાંગમાં
    $\lforall[y][!C_1'(y)]$
    ઉમેરવા શિથિલીકરણની સ્વીકાર્યતા
    વાપરીએ છીએ).
    પહેલી !!{proof}માં
    \RightR{\lforall}ની સ્વચલ-શરત
    પૂરી થાય છે: કારણ કે
    $b \notin \Lang L(\Gamma_1, \Delta_1,
    \lexists[x][!A(x)])$ અને
    $!C_1'$ની વ્યાખ્યા એવી છે કે
    તેમાં~$b$ નથી.
    
    બીજા કિસ્સામાં આગમન-પૂર્વધારણા
    અને~$\lexists[y][!C_1'(y)]$
    ઉમેરતા શિથિલીકરણની સ્વીકાર્યતા વડે:
    \[
    \AxiomC{}
    \Deduce$\Gamma_1 \fCenter
      \Delta_1, \lexists[x][!A(x)], !A(b), \lexists[y][!C_1'(y)], !C_1'(b)$
    \RightLabel{\RightR{\lexists}}
    \UnaryInf$\Gamma_1 \fCenter
      \Delta_1, \lexists[x][!A(x)], \lexists[y][!C_1'(y)], !C_1'(b)$
    \RightLabel{\RightR{\lexists}}
    \UnaryInf$\Gamma_1 \fCenter
      \Delta_1, \lexists[x][!A(x)], \lexists[y][!C_1'(y)]$
    \DisplayProof
    \]
    અને
    \[
    \AxiomC{}
    \Deduce$!C_1'(b), \Gamma_2 \fCenter \Delta_2$
    \RightLabel{\LeftR{\lexists}}
    \UnaryInf$\lexists[y][!C_1'(y)], \Gamma_2 \fCenter \Delta_2$
    \DisplayProof
    \]
    બીજી !!{proof}માં
    \LeftR{\lexists}ની સ્વચલ-શરત
    પૂરી થાય છે, કારણ કે
    $b \notin \Lang L(\Gamma_2, \Delta_2)$
    અને~$!C_1'$ની વ્યાખ્યા એવી છે
    કે તેમાં~$b$ નથી.
  \end{enumerate}
  બીજા કિસ્સા સમાન છે;
  તેમને કસરત તરીકે છોડીએ છીએ.
\end{proof}

\begin{prob}
માનો કે~$\lnand$ સંયોજક માટે
આ નિયમો હોય:
\[
\Axiom$\Gamma \fCenter \Delta, !A$
\Axiom$\Gamma \fCenter \Delta, !B$
\RightLabel{\LeftR{\lnand}}
\BinaryInf$!A \lnand !B, \Gamma \fCenter \Delta$
\DisplayProof
\quad
\Axiom$!A, !B, \Gamma \fCenter \Delta$
\RightLabel{\RightR{\lnand}}
\UnaryInf$\Gamma \fCenter \Delta, !A \lnand !B$
\DisplayProof
\]
\olref{lem:maehara}ની સાબિતીમાં
જ્યાં~$\pi$નો અંત આમાંથી કોઈ
નિયમથી થાય છે તે કિસ્સા કરીને
માએહારાનું ઉપપ્રમેય હજી સાચું છે
એ બતાવો.
\end{prob}

માએહારાનું ઉપપ્રમેય મળ્યું છે;
હવે તે વડે ક્રેગનું અંતર્વેશન
પ્રમેય સાબિત કરીએ.

\begin{proof}[\olref{thm:interpolation}ની સાબિતી]
  માનો કે $!A \Sequent !B$
  !!{provable} છે.
  $\maeh{!A}{\ } \Sequent \maeh{\ }{!B}$
  પર \olref{lem:maehara} લગાવતાં
  અંતર્વેશક~$!C$ મળે છે.
  આપણને
  $\Proves !A \Sequent !C$ અને
  $\Proves !C \Sequent !B$ મળે છે.
\end{proof}

માનો કે~$\Gamma$ એ~$n$-સ્થાની
વિધેય~$R$ ધરાવતી
!!a{language}~$\Lang L$નો
સિદ્ધાંત (!!{sentence}sનો ગણ) છે.
~$\Gamma$ એ~$R$ને
\emph{ગૂઢ રીતે વ્યાખ્યાયિત કરે છે}
ત્યારે અને માત્ર ત્યારે જ જ્યારે
\begin{align*}
\Gamma, \Gamma'
& \Entails \lforall[x_1][\dots\lforall[x_n][(R(x_1,\dots,x_n) \liff
R'(x_1, \dots, x_n))]],
\intertext{અહીં $\Gamma'$ એ $\Gamma$માં $R$ના બધા દૃષ્ટાંતોને $R' \notin \Lang L$થી બદલીને મળ્યું છે. $\Gamma$ એ $R$ને \emph{સ્પષ્ટ રીતે વ્યાખ્યાયિત કરે છે} ત્યારે અને માત્ર ત્યારે જ જ્યારે $R$ ન ધરાવતું એવું $!A(x_1, \dots, x_n)$ છે કે}
\Gamma &\Entails
\lforall[x_1][\dots\lforall[x_n][(R(x_1,\dots,x_n) \liff !A(x_1,
\dots, x_n))]].
\end{align*}
સ્પષ્ટ છે કે જો~$\Gamma$ એ~$R$ને
સ્પષ્ટ રીતે વ્યાખ્યાયિત કરે,
તો ગૂઢ રીતે પણ કરે છે.
~$R$ માટે વ્યસ્ત દિશા તરત સ્પષ્ટ નથી,
પણ તે પણ સાચી છે:

\begin{lem}[બેથનું વ્યાખ્યેયતા પ્રમેય]
  જો~$\Gamma$ એ~$R$ને ગૂઢ રીતે
  વ્યાખ્યાયિત કરે, તો સ્પષ્ટ
  રીતે પણ વ્યાખ્યાયિત કરે છે.
\end{lem}

\begin{proof}
  માનો કે~$\Gamma$ એ~$R$ને ગૂઢ રીતે
  વ્યાખ્યાયિત કરે છે.
  ઉપર મુજબ~$R'$ અને~$\Gamma'$
  લો તથા
  $\Lang L' = \Lang L(\Gamma') =
  \Lang L \setminus \{R\}
  \cup \{R'\}$.
  પૂર્વધારણાથી મળે છે:
  \begin{align*}
    \Gamma, \Gamma' & 
    \Sequent \lforall[x_1][\dots\lforall[x_n][(R(x_1,\dots,x_n) 
    \liff R'(x_1, \dots, x_n))]]
  \intertext{$c_1$, \dots,~$c_n$ એ $\Gamma$માં ન આવતાં !!{constant}s હો. ઉલટાવી શકવાના ઉપપ્રમેયથી મળે છે:}
    \Gamma, \Gamma', R(c_1,\dots,c_n) & \Sequent  R'(c_1, \dots, c_n)
  \intertext{માએહારાનું ઉપપ્રમેય આના પર લગાવો:}
    \maeh{\Gamma, R(c_1,\dots,c_n)}{\Gamma'} & \Sequent
      \maeh{\ }{R'(c_1, \dots, c_n)}
  \intertext{ત્યારે એવું $!A \ident !A(x_1, \dots, x_n)$ મળે છે કે}
    \Gamma, R(c_1,\dots,c_n) & \Sequent !A(c_1, \dots, c_n) \\
    !A(c_1, \dots, c_n), \Gamma' & \Sequent R'(c_1, \dots, c_n)
  \intertext{આ બંનેની !!{proof}s છે. માએહારાના ઉપપ્રમેય મુજબ $\Lang L(!A) \subseteq \Lang L(\Gamma) \setminus \{R\}$; એટલે $!A$માં $R$ નથી. બીજા સિક્વન્ટની !!{proof}માં $R'$ને સર્વત્ર $R$થી બદલો; ત્યારે આની !!a{proof} મળે છે:}
    !A(c_1, \dots, c_n), \Gamma & \Sequent R(c_1, \dots, c_n)
  \end{align*}
  \sourcecorrection{OLMD-207}{મૂળમાં અંતર્વેશકની ભાષા સામાન્ય ભાષાના બરાબર છે અને અપરિભાષિત $L_1,L_2$ વાપરે છે; ઉપપ્રમેય માત્ર સમાવેશ આપે છે, જે અહીં પૂરતો છે.}
  બંને દિશાના અનુસરણો માટે
  \RightR{\lif}, પછી દ્વિમુખી
  સંબંધની રજૂઆત અને
  \RightR{\lforall} લગાવી
  આની !!a{proof} મેળવો:
  \sourcecorrection{OLMD-208}{મૂળમાં માત્ર અનુસરણ-જમણો અને સર્વમાન્યક-જમણો નિયમ લખેલા છે; બે દિશા ભેગી કરીને દ્વિમુખી સંબંધ બનાવવાનું મધ્યવર્તી પગલું પણ જરૂરી છે.}
  \[
    \Gamma \Sequent \lforall[x_1][\dots\lforall[x_n][(R(x_1,\dots,x_n) \liff !A(x_1, \dots, x_n))]].
  \]
  આ બતાવે છે કે
  $!A(x_1, \dots, x_n)$
  એ~$\Gamma$માં~$R$ને
  સ્પષ્ટ રીતે વ્યાખ્યાયિત કરે છે.
  \sourcecorrection{OLMD-209}{મૂળ અંતિમ વાક્યમાં $!R$ લખાયું છે, જ્યારે $R$ અહીં વ્યાખ્યાયિત થતું વિધેયચિહ્ન છે.}
\end{proof}

અંતર્વેશનને સમકક્ષ ત્રીજું પરિણામ
(જેને પણ માએહારાના ઉપપ્રમેયથી
સાબિત કરી શકાય) આ છે:
જો બે સુસંગત સિદ્ધાંતો~$\Gamma_1$
અને~$\Gamma_2$માં તેમના સામાન્ય
ભાષાનો પૂર્ણ સિદ્ધાંત~$\Gamma$
સમાયેલો હોય, એટલે
$\Lang L(\Gamma) =
\Lang L(\Gamma_1) \cap \Lang L(\Gamma_2)$,
તો $\Gamma_1 \cup \Gamma_2$
સુસંગત છે.
\sourcecorrection{OLMD-210}{મૂળમાં $\Lang L(\Gamma)$ને માત્ર સામાન્ય ભાષાનો ઉપગણ કહે છે; આગળની પૂર્ણતાની દલીલ દરેક સામાન્ય-ભાષાના વાક્ય પર ચાલે છે, તેથી અહીં ભાષાઓની સમાનતા જરૂરી છે.}

યાદ કરો: પૂર્ણ સિદ્ધાંત એવો છે કે તેની
ભાષાના દરેક !!{sentence}~$!A$ માટે
કાં તો $\Gamma \Proves !A$ અથવા
$\Gamma \Proves \lnot !A$.
કારણ કે~$\Gamma_1$ અને~$\Gamma_2$
એ~$\Gamma$ના સુસંગત વિસ્તારો છે,
$\Gamma$ પણ સુસંગત છે.
એથી જો~$!A$ એ
$\Lang L(\Gamma_1) \cap
\Lang L(\Gamma_2)$ની
!!a{sentence} હોય, તો
$\Gamma_1 \Entails !A$ અને
$\Gamma_2 \Entails \lnot !A$
બંને હોઈ શકે નહીં.
વિરોધ માટે માનો કે એવું~$!A$ છે.
$\Gamma_1 \Entails !A$ છે;
આથી પૂર્ણતા મુજબ
$\Gamma \Entails !A$ જ હોવું જોઈએ.
\sourcecorrection{OLMD-211}{મૂળમાં $\Gamma_1\Entails !A$ પરથી સીધું $\Gamma\Entails !A$ કાઢ્યું છે; નાનું સિદ્ધાંત મોટાના બધા પરિણામો આપતું નથી. પૂર્ણતા અને $\Gamma_1$ની સુસંગતતા વડે જ આ નિષ્કર્ષ મળે છે.}
નહીં તો~$\Gamma$ની પૂર્ણતાથી
$\Gamma \Entails \lnot !A$,
અને $\Gamma \subseteq \Gamma_1$
હોવાથી
$\Gamma_1 \Entails \lnot !A$ પણ,
જે~$\Gamma_1$ને અસુસંગત બનાવે.
તેથી $\Gamma \Entails !A$;
અને $\Gamma \subseteq \Gamma_2$
હોવાથી
$\Gamma_2 \Entails !A$.
પણ પૂર્વધારણા મુજબ
$\Gamma_2 \Entails \lnot !A$ પણ છે,
જે~$\Gamma_2$ની સુસંગતતા સામે છે.
\sourcecorrection{OLMD-212}{મૂળ અંતિમ સૂત્રમાં $\Gamma\Entails/\lnot !A$ લખ્યું છે; વિરોધ ખરેખર $\Gamma_2$માં $!A$ અને $\lnot !A$ બંને ફલિત થવાથી મળે છે.}

સાબિતી માટે એટલું જ જરૂરી છે કે
સામાન્ય
!!{language}~$\Lang L(\Gamma_1)
\cap \Lang L(\Gamma_2)$માં
એવું કોઈ !!a{sentence}~$!A$ ન હોય
જેના માટે
$\Gamma_1 \Entails !A$ અને
$\Gamma_2 \Entails \lnot !A$.
હવે માએહારાના ઉપપ્રમેય વડે
સાબિતી-સૈદ્ધાંતિક સાબિતી આપીએ.

\begin{thm}[રૉબિન્સનનું સંયુક્ત સુસંગતતા પ્રમેય]
માનો કે~$\Gamma_1$ અને~$\Gamma_2$
સુસંગત સિદ્ધાંતો છે, અને
!!{language}
$\Lang L(\Gamma_1) \cap
\Lang L(\Gamma_2)$માં એવું
કોઈ~$!A$ નથી કે
$\Gamma_1 \Entails !A$ અને
$\Gamma_2 \Entails \lnot !A$.
તો~$\Gamma_1 \cup \Gamma_2$
સુસંગત છે.
\end{thm}

\begin{proof}
  વિરોધ માટે માનો કે
  $\Gamma_1 \cup \Gamma_2$
  અસુસંગત છે. ત્યારે
  $\Gamma_1, \Gamma_2 \Sequent \ $
  સિક્વન્ટની !!a{proof} છે.
  $\maeh{\Gamma_1}{\Gamma_2}
  \Sequent \maeh{\ }{\ }$
  પર માએહારાનું ઉપપ્રમેય
  લગાવતાં
  !!{language}~$\Lang L(\Gamma_1)
  \cap \Lang L(\Gamma_2)$માં
  એવી !!a{sentence}~$!A$ મળે છે કે
  $\Proves \Gamma_1 \Sequent !A$
  અને
  $\Proves !A, \Gamma_2 \Sequent \quad$.
  બીજા સિક્વન્ટથી
  $\Proves \Gamma_2 \Sequent \lnot !A$
  મળે છે. પરંતુ એવું કોઈ
  !!{sentence} નથી એમ ધારેલું છે.
\end{proof}

ઉપર સિદ્ધાંતો~$\Gamma$,
$\Gamma_1$ અને~$\Gamma_2$
સાન્ત છે એવું ગર્ભિત રીતે
ધાર્યું છે. સઘનતાથી પરિણામો
અનંત સિદ્ધાંતો માટે પણ સાચાં છે.

\end{document}
